% !TeX root = main.tex
\section{Flow Dynamics Conventions}

\subsection{Current Simple-Model Direction}

Scheduled particle birth, lifecycle gating through \texttt{Data.w}, and
outlet retirement have been retired from the live Python
\texttt{ForceDynamics} path.  Mobile particles are generated active from
frame zero.  The current verification target is the simple particle model:
mobile particles, boundary-particle locality markers, and stationary
parametric walls generated from \texttt{curve\_wall\_segments}.  Reservoir,
piston, and CD-nozzle flow remain separate later efforts until the simple
Python and GLSL paths agree.

The former deferred-birth design remains documented as a historical
investigation in \texttt{ProblemEvolution.tex}.  It is not an active
generation or dynamics rule.

The active generic generator contract is:
\begin{itemize}
    \item mobile particles are active at frame zero;
    \item boundary particles use \texttt{ptype=1};
    \item \texttt{Data.z} is not used for wall evaluator selection;
    \item all stationary wall geometry is in
          \texttt{curve\_wall\_segments};
    \item \texttt{death\_bounds} define permanent mobile-particle death;
    \item \texttt{cell\_array\_width},
          \texttt{cell\_array\_height}, and
          \texttt{cell\_array\_depth} define the rectangular cell domain;
    \item legacy derived \texttt{boundary\_*} test-file fields are not
          emitted by \texttt{GenericGenData}.
\end{itemize}

\subsection{Horizontal Reservoir Generator}

The active replacement work begins in:
\begin{verbatim}
gbase/BoundaryParticleReservoirHorizontal.py
cfg_gendata/BoundaryParticleReservoirHorizontal.cfg
\end{verbatim}

The generator owns the horizontal orientation.  Its cfg does not select a
wall model, wall function, conduit axis, birth schedule, or retirement plane.
The reduced Python profile remains selected through
\texttt{sub\_BPRPPy.cfg}; the Vulkan profile is deferred until Python
geometry and dynamics are verified.

The packed physical domain is configured as:
\begin{verbatim}
chamber_bounds = [
    x_start, x_end,
    y_bottom, y_top,
    z_front, z_back
];
\end{verbatim}

For the horizontal model, positive (x) is axial.  Every chamber bound must
lie inside the explicit rectangular cell domain defined by
\texttt{cell\_array\_width}, \texttt{cell\_array\_height}, and
\texttt{cell\_array\_depth}.

The current two-dimensional convention is:
\[
z_{\mathrm{front}}=0,\qquad
z_{\mathrm{back}}=3,\qquad
z_{\mathrm{particle}}=1.
\]
The cell-array depth must therefore be greater than three.

\subsection{Reservoir Packing}

The current generator calculates reservoir capacity and emits the complete
centered mobile-particle grid.  With particle radius (r) and configured
surface-to-surface separation (s), center spacing is:
\[
\ell=2r+s.
\]

Initial particle centers reserve the wall-contact clearance:
\[
c=r(1+o)+\epsilon,
\]
where (o) is \texttt{wall\_contact\_offset}.  The usable center interval
on each packed axis is the physical interval reduced by (c) at both ends.
For usable span (L), the count is:
\[
N=\left\lfloor\frac{L}{\ell}\right\rfloor+1.
\]
The occupied grid is centered in any remaining space so opposing wall margins
are equal.

Generated mobile particles are active from frame zero.  Piston motion uses
the separately configured \texttt{piston\_velocity}.

\subsection{Boundary Marker Geometry}

The earlier reservoir prototype used two named finite-line collections:
\begin{verbatim}
linear_chamber_segments = (...);
linear_wall_segments = (...);
\end{verbatim}
Each row is \texttt{[A,B,C,xStart,yStart,xEnd,yEnd,wallFlag]} and represents
the finite line \(Ax+By+C=0\).  The normalized vector \((A,B)\) is the outward
wall normal.  The collections distinguish chamber walls from attached
conduit or nozzle walls without introducing different collision equations.

The simple generic path replaces those line-specific arrays with
\texttt{curve\_wall\_segments}.  Straight walls are represented as cubic
segments with zero higher-order coefficients, so horizontal, vertical,
sloped, and curved walls share one evaluator.  Markers are no farther than one
cell apart and shared coordinates are emitted once per physical wall flag.
Boundary markers remain occupancy placeholders rather than mobile dynamics
sources.

The fixed vertical inlet marker row has been removed.  The analytic piston
occupies that physical role without becoming a boundary particle.  The
generator writes both binary particle data and test metadata.

\subsection{Current Piston Work}

\texttt{piston\_start\_frame} controls when the rear plane begins moving.
The piston starts at \texttt{chamber\_bounds.x\_start}, advances with
\texttt{piston\_velocity}, and stops at
\texttt{chamber\_bounds.x\_end}.  Each
active source evaluates it once before target scanning.  Python diagnostics
report piston position, velocity, contact count, impulse, and momentum
transfer.  This work is deferred while \texttt{TestPythonBoundarySpheres}
becomes the simple low-particle Vulkan parity target.  When piston and
reservoir work resumes, it should inherit the same generic parametric wall
contract instead of reintroducing finite-line or evaluator-ID branches.
